\scanpage{8}
% Typescript upside down in the archive scan.
\mspage{- 19 -}
n\textsuperscript{o} 391

terminer la démonstration, soient $y$ un point de $C$, $z$ un point de $C'=\bigcap\limits_{H\in\underline{\unsure{M|H}}}D_H(C)$. Comme $\underline{H}$ est localement fini, le segment fermé d'extrémités $y$, $z$ ne peut rencontrer qu'un nombre fini d'hyperplans de $\underline{H}$ ; il n'y a par suite qu'un nombre fini d'hyperplans $H_1,..,H_q$ de $\underline{H}$ tels que $z\notin D_{H_i}(C)$ pour $1\leq i\leq q$. Par définition de $C'$, les $H_i$ ne sont pas des murs de $C$, donc, par récurrence sur $i$, on voit que $C$ est l'intersection des $D_H(C)$ pour les $H\in\underline{H}$ \uline{distincts des} $H_i$. Cela prouve que $z\in C$, donc $C'\subset C$, et par suite $C'=C$.

COROLLAIRE 2.- \uline{Soit} $\underline{H}$ \uline{un ensemble localement fini d'hyperplans affines}. \uline{Alors tout hyperplan} $H\in\underline{H}$ \uline{est mur mitoyen de deux chambres déterminées par} $\underline{H}$.

Il suffit d'appliquer la prop. 10 à l'ensemble $\underline{H}-\{H\}$.

COROLLAIRE 3.- \uline{Soient} $\underline{H}$ \uline{un ensemble localement fini d'hyperplans affines,} $C$ \uline{une chambre déterminée par} $\underline{H}$, $H_{\unsure{o|0}}$ \uline{un mur de} $C$. \uline{L'intérieur de} $\bar{C}\cap H_o$ \uline{dans} $H_o$ \uline{est une chambre dans} $H_o$, \uline{formée des points de} $\bar{C}\cap H_o$ \uline{qui n'appartiennent à aucun hyperplan} $H\ne H_o$ \uline{de} $\underline{H}$ ; \uline{ce sont aussi les points} $x$ \uline{de} $\bar{C}\cap H_o$ \uline{qui sont adhérents à} $C$ \uline{et à une seule des autres chambres déterminées par} $\underline{H}$, \uline{et cette chambre} $C_1$ \uline{est mitoyenne à} $C$ \uline{et est la même pour tous les points} $x$ \uline{intérieurs à} $\bar{C}\cap H_o$.

Le fait que $H_o$ soit un mur de $C$ entraîne que $C_o=\bigcap\limits_{H\in\underline{H},H\ne H_o}D_H(C)$ est distincte de $C$ (cor. 1) ; $C$ est donc l'une des deux chambres notées $C'$, $C''$ dans l'énoncé de la prop. 10 ; supposons par exemple que $C=C'$ ; avec les mêmes notations, l'intérieur de $\bar{C}\cap H_o$ dans $H_o$ est $\Gamma$, et on a vu dans la prop. 10 que les points de $\Gamma$ ne sont adhérents qu'à deux des chambres déterminées par $\underline{H}$, $C'$ et la chambre mitoyenne $C''$. Au contraire, si un point $x\in\bar{C}\cap H_o=\bar{\Gamma}$ n'appartient pas à $\Gamma$, il appartient à un des hyperplans $H\ne H_o$ de $\underline{H}$ ; il est
% archive sheet number: 7
